\ifdefined\gkver\else\def\gkver{student}\fi
\documentclass[\gkver]{gaokaozhenti}
\begin{document}

\chapter{2023年1月浙江}

\begin{enumerate}
\item
%% number: 1
%% paperName: 2023年1月浙江选考第 1 题 3 分
%% typeId: 1
%% type: 单选题
%% chapter: 运动和力
%% point: 力学单位制
%% method: 无
%% score: 3
%% degree: 950
%% duplicateId: 0
%% body:
下列属于国际单位制中基本单位符号的是\xzanswer{B}
\fourchoices[answer=B]
{J}
{K}
{W}
{Wb}
%% answer: B
%% memo:
\memoanswer{
国际单位制中的七个基本单位是千克、米、秒、安培、开尔文、坎德拉、摩尔，符号分别是 kg、m、s、A、K、cd、mol，其余单位都属于导出单位。

故选 B。
}

\item
%% number: 2
%% paperName: 2023年1月浙江选考第 2 题 3 分
%% typeId: 1
%% type: 单选题
%% chapter: 力物体的平衡
%% point: 三力平衡
%% method: 无
%% score: 3
%% degree: 850
%% duplicateId: 0
%% body:
如图所示，轻质网兜兜住重力为 $G$ 的足球，用轻绳挂于光滑竖直墙壁上的 A 点，轻绳的拉力为 $F_{T}$，墙壁对足球的支持力为 $F_{N}$，则\xzanswer{C}
\onepicture{figs/02.png}
\fourchoices[answer=C]
{$F_{T} < F_{N}$}
{$F_{T} = F_{N}$}
{$F_{T} > G$}
{$F_{T} = G$}
%% answer: C
%% memo:
\memoanswer{
对网兜和足球受力分析，设轻绳与竖直墙面夹角为 $\theta$，由平衡条件

 $F_{T}=\frac{G}{\cos\theta}=\frac{F_{N}}{\sin\theta},\quad F_{N}=G\tan\theta$

可知  $F_{T} > G, F_{T} > F_{N}$

故选 C。
}

\item
%% number: 3
%% paperName: 2023年1月浙江选考第 3 题 3 分
%% typeId: 1
%% type: 单选题
%% chapter: 直线运动
%% point: 参考系
%% method: 无
%% score: 3
%% degree: 950
%% duplicateId: 0
%% body:
“神舟十五号”飞船和空间站“天和”核心舱成功对接后，在轨运行如图所示，则\xzanswer{C}
\onepicture{figs/03.jpg}
\fourchoices[answer=C]
{选地球为参考系，“天和”是静止的}
{选地球为参考系，“神舟十五号”是静止的}
{选“天和”为参考系，“神舟十五号”是静止的}
{选“神舟十五号”为参考系，“天和”是运动的}
%% answer: C
%% memo:
\memoanswer{
AB．“神舟十五号”飞船和空间站“天和”核心舱成功对接后，在轨绕地球做圆周运动，选地球为参考系，二者都是运动的，AB 错误；

CD．“神舟十五号”飞船和空间站“天和”核心舱成功对接后，二者相对静止，C 正确，D 错误。

故选 C。
}

\item
%% number: 4
%% paperName: 2023年1月浙江选考第 4 题 3 分
%% typeId: 1
%% type: 单选题
%% chapter: 功和能
%% point: 功能原理
%% method: 无
%% score: 3
%% degree: 650
%% duplicateId: 0
%% body:
一位游客正在体验蹦极，绑上蹦极专用的橡皮绳后从跳台纵身而下。游客从跳台下落直到最低点过程中\xzanswer{B}
\fourchoices[answer=B]
{弹性势能减小}
{重力势能减小}
{机械能保持不变}
{绳一绷紧动能就开始减小}
%% answer: B
%% memo:
\memoanswer{
游客从跳台下落，开始阶段橡皮绳未拉直，只受重力作用做匀加速运动，下落到一定高度时橡皮绳开始绷紧，游客受重力和向上的弹力作用，弹力从零逐渐增大，游客所受合力先向下减小后向上增大，速度先增大后减小，到最低点时速度减小到零，弹力达到最大值。

A．橡皮绳绷紧后弹性势能一直增大，A 错误；

B．游客高度一直降低，重力一直做正功，重力势能一直减小，B 正确；

C．下落阶段橡皮绳对游客做负功，游客机械能减少，转化为弹性势能，C 错误；

D．绳刚绷紧开始一段时间内，弹力小于重力，合力向下做正功，游客动能在增加；当弹力大于重力后，合力向上对游客做负功，游客动能逐渐减小，D 错误。

故选 B。
}

\item
%% number: 5
%% paperName: 2023年1月浙江选考第 5 题 3 分
%% typeId: 1
%% type: 单选题
%% chapter: 曲线运动
%% point: 斜抛运动
%% method: 无
%% score: 3
%% degree: 550
%% duplicateId: 0
%% body:
如图所示，在考虑空气阻力的情况下，一小石子从 O 点抛出沿轨迹 OPQ 运动，其中 P 是最高点。若空气阻力大小与瞬时速度大小成正比，则小石子竖直方向分运动的加速度大小\xzanswer{A}
\onepicture{figs/05.png}
\fourchoices[answer=A]
{O 点最大}
{P 点最大}
{Q 点最大}
{整个运动过程保持不变}
%% answer: A
%% memo:
\memoanswer{
由于空气阻力大小与瞬时速度大小成正比，小石子在 O 点时速度斜向上方，此时速度最大，空气阻力斜向下方最大，上升过程与竖直方向夹角最小，故此时空气阻力分解在竖直方向最大，根据牛顿第二定律可知此时竖直方向分运动的加速度最大。

故选 A。
}

\item
%% number: 6
%% paperName: 2023年1月浙江选考第 6 题 3 分
%% typeId: 1
%% type: 单选题
%% chapter: 机械波
%% point: 干涉衍射
%% method: 无
%% score: 3
%% degree: 650
%% duplicateId: 0
%% body:
主动降噪耳机能收集周围环境中的噪声信号，并产生相应的抵消声波，某一噪声信号传到耳膜的振动图像如图所示，取得最好降噪效果的抵消声波（声音在空气中的传播速度为 340 $\Ums$）\xzanswer{B}
\onepicture{figs/06.png}
\fourchoices[answer=B]
{振幅为 2$A$}
{频率为 100 $\UHz$}
{波长应为 1.7 $\Um$ 的奇数倍}
{在耳膜中产生的振动与图中所示的振动同相}
%% answer: B
%% memo:
\memoanswer{
主动降噪耳机是根据波的干涉条件，抵消声波与噪声的振幅、频率相同，相位相反，叠加后才能相互抵消来实现降噪的。

A．抵消声波与噪声的振幅相同，也为 $A$，A 错误；

B．抵消声波与噪声的频率相同，由 $f = \frac{1}{T}$ = 100 $\UHz$。B 正确；

C．抵消声波与噪声的波速、频率相同，则波长也相同，为 $\lambda = vT$ = 340$\times$0.01 $\Um$ = 3.4 $\Um$。C 错误；

D．抵消声波在耳膜中产生的振动与图中所示的振动反相，D 错误。

故选 B。
}

\item
%% number: 7
%% paperName: 2023年1月浙江选考第 7 题 3 分
%% typeId: 1
%% type: 单选题
%% chapter: 电磁感应
%% point: 楞次定律推广
%% method: 无
%% score: 3
%% degree: 550
%% duplicateId: 0
%% body:
如图甲所示，一导体杆用两条等长细导线悬挂于水平轴 OO$'$，接入电阻 $R$ 构成回路．导体杆处于竖直向上的匀强磁场中，将导体杆从竖直位置拉开小角度由静止释放，导体杆开始下摆。当 $R = R_{0}$ 时，导体杆振动图像如图乙所示。若横纵坐标皆采用图乙标度，则当 $R$ = 2$R_{0}$ 时，导体杆振动图像是\xzanswer{B}
\twopicture[labelA=20231月浙江07a, labelB=20231月浙江07b]
{figs/07a.png}
{\includesvg{figs/07b}}
%% answer: B
%% memo:
\memoanswer{
导体杆切割磁感线时，回路中产生感应电流，由楞次定律可得，导体杆受到的安培力总是阻碍导体棒的运动。当 $R$ 从 $R_{0}$ 变为 2$R_{0}$ 时，回路中的电阻增大，则电流减小，导体杆所受安培力减小，即导体杆在摆动时所受的阻力减弱，所杆从开始摆动到停止，运动的路程和经历的时间变长。

故选 B。
}

\item
%% number: 8
%% paperName: 2023年1月浙江选考第 8 题 3 分
%% typeId: 1
%% type: 单选题
%% chapter: 磁场
%% point: 带电粒子在磁场中的运动
%% method: 无
%% score: 3
%% degree: 650
%% duplicateId: 0
%% body:
某兴趣小组设计的测量大电流的装置如图所示，通有电流 $I$ 的螺绕环在霍尔元件处产生的磁场 $B = k_{1}I$，通有待测电流 $I'$ 的直导线 ab 垂直穿过螺绕环中心，在霍尔元件处产生的磁场 $B' = k_{2}I'$。调节电阻 $R$，当电流表示数为 $I_{0}$ 时，元件输出霍尔电压 $U_{H}$ 为零，则待测电流 $I'$ 的方向和大小分别为\xzanswer{D}
\onepicture{figs/08.png}
\fourchoices[answer=D]
{a$\to$b，$\frac{{{k_2}}}{{{k_1}}} I_{0}$}
{a$\to$b，$\frac{{{k_1}}}{{{k_2}}}I_{0}$}
{b$\to$a，$\frac{{{k_2}}}{{{k_1}}} I_{0}$}
{b$\to$a，$\frac{{{k_1}}}{{{k_2}}}I_{0}$}
%% answer: D
%% memo:
\memoanswer{
根据安培定则可知螺绕环在霍尔元件处产生的磁场方向向下，则要使元件输出霍尔电压 $U_{H}$ 为零，直导线 ab 在霍尔元件处产生的磁场方向应向上，根据安培定则可知待测电流 $I'$ 的方向应该是 b$\to$a；元件输出霍尔电压 $U_{H}$ 为零，则霍尔元件处合场强为 0，所以有

$k_{1}I_{0} = k_{2}I'$

解得 $I' = \frac{{{k_1}}}{{{k_2}}} I_{0}$

故选 D。
}

\item
%% number: 9
%% paperName: 2023年1月浙江选考第 9 题 3 分
%% typeId: 1
%% type: 单选题
%% chapter: 原子和原子核
%% point: 半衰期
%% method: 无
%% score: 3
%% degree: 650
%% duplicateId: 0
%% body:
宇宙射线进入地球大气层与大气作用会产生中子，中子与大气中的氮 14 会产生以下核反应：\ce{^{14}_{7}N} + \ce{^{1}_{0}n} $\to$ \ce{^{14}_{6}C} + \ce{^{1}_{1}H}，产生的 \ce{^{14}_{6}C} 能自发进行 $\beta$ 衰变，其半衰期为 5 730 年，利用碳 14 的衰变规律可推断古木的年代．下列说法正确的是\xzanswer{D}
\fourchoices[answer=D]
{\ce{^{14}_{6}C} 发生 $\beta$ 衰变的产物是 \ce{^{15}_{7}N}}
{$\beta$ 衰变辐射出的电子来自于碳原子的核外电子}
{近年来由于地球的温室效应，引起 \ce{^{14}_{6}C} 的半衰期发生微小变化}
{若测得一古木样品的 \ce{^{14}_{6}C} 含量为活体植物的 $\frac{1}{4}$，则该古木距今约为 11 460 年}
%% answer: D
%% memo:
\memoanswer{
A．根据 \ce{^{14}_{6}C} $\to$ \ce{^{14}_{7}N} + \ce{^{0}_{-1}e}

即 \ce{^{14}_{6}C} 发生 $\beta$ 衰变的产物是 \ce{^{14}_{7}N}，选项 A 错误；

B．$\beta$ 衰变辐射出的电子来自于原子核内的中子转化为质子时放出的电子，选项 B 错误；

C．半衰期是核反应，与外界环境无关，选项 C 错误；

D．若测得一古木样品的 \ce{^{14}_{6}C} 含量为活体植物的 $\frac{1}{4}$，可知经过了2个半衰期，则该古木距今约为 5 730$\times$2年=11 460年，选项 D 正确。

故选 D。
}

\item
%% number: 10
%% paperName: 2023年1月浙江选考第 10 题 3 分
%% typeId: 1
%% type: 单选题
%% chapter: 曲线运动
%% point: 天体运动
%% method: 无
%% score: 3
%% degree: 550
%% duplicateId: 0
%% body:
太阳系各行星几平在同一平面内沿同一方向绕太阳做圆周运动．当地球恰好运行到某地外行星和太阳之间，且三者几乎排成一条直线的现象，称为“行星冲日”，已知地球及各地外行星绕太阳运动的轨道半径如下表：

\begin{center}
\begin{tabular}{|c|}
\hline
行星名称 \\ \hline
\end{tabular}
\end{center}

\begin{center}
\begin{tblr}{|c|c|c|c|c|c|c|}
\hline
 & 地球 & 火星 & 木星 & 土星 & 天王星 & 海王星 \\
\hline
$R/\mathrm{AU}$ & 1.0 & 1.5 & 5.2 & 9.5 & 19 & 30 \\
\hline
\end{tblr}
\end{center}

则相邻两次“冲日”时间间隔约为\xzanswer{B}
\fourchoices[answer=B]
{火星 365 天}
{火星 800 天}
{天王星 365 天}
{天王星 800 天}
%% answer: B
%% memo:
\memoanswer{
根据开普勒第三定有
$$ \frac{T^{2}}{R^{3}}=\frac{T_{ \text{地} }^{2}}{R_{ \text{地} }^{3}} $$

解得

$$ T=\sqrt{(\frac{R}{R_{ \text{地} }})^{3}T_{ \text{地} }} $$

设相邻两次“冲日”时间间隔为 t，则

 $2\pi=\left(\frac{2\pi}{T_{\text{地}}}-\frac{2\pi}{T}\right)t$

解得

 $t=\frac{TT_{\text{地}}}{T-T_{\text{地}}}=\frac{T_{\text{地}}}{1-\sqrt{\frac{R_{\text{地}}^{3}}{R^{3}}}}$

由表格中的数据可得

 $t_{\text{火}}=\frac{T_{\text{地}}}{1-\sqrt{\frac{R_{\text{地}}^{3}}{R_{\text{火}}^{3}}}}\approx800$天

 $t_{\text{天}}=\frac{T_{\text{地}}}{1-\sqrt{\frac{R_{\text{地}}^{3}}{R_{\text{天}}^{3}}}}\approx369$天

故选 B。
}

\item
%% number: 11
%% paperName: 2023年1月浙江选考第 11 题 3 分
%% typeId: 1
%% type: 单选题
%% chapter: 光学
%% point: 光子说
%% method: 建立模型
%% score: 3
%% degree: 550
%% duplicateId: 0
%% body:
被誉为“中国天眼”的大口径球面射电望远镜已发现 660 余颗新脉冲星，领先世界。天眼对距地球为 $L$ 的天体进行观测，其接收光子的横截面半径为 $R$。若天体射向天眼的辐射光子中，有 $\eta$（$\eta$ < 1）倍被天眼接收，天眼每秒接收到该天体发出的频率为 $\nu$ 的 $N$ 个光子。普朗克常量为 $h$，则该天体发射频率为 $\nu$ 光子的功率为\xzanswer{A}
\onepicture{figs/11.jpg}
\fourchoices[answer=A]
{$\frac{{4N{L^2}h\nu }}{{{R^2}\eta }}$}
{$\frac{{2N{L^2}h\nu }}{{{R^2}\eta }}$}
{$\frac{{\eta {L^2}h\nu }}{{4{R^2}N}}$}
{$\frac{{\eta {L^2}h\nu }}{{2{R^2}N}}$}
%% answer: A
%% memo:
\memoanswer{
设天体发射频率为 v 光子的功率为 P，由题意可知

 $Pt \times \frac{\pi R^{2}}{4\pi L^{2}} \times \eta = Nhvt$

解得  $P = \frac{4NL^{2}hv}{R^{2}\eta}$

故选 A。
}

\item
%% number: 12
%% paperName: 2023年1月浙江选考第 12 题 3 分
%% typeId: 6
%% type: 计算题
%% chapter: 电场
%% point: 带电粒子的偏转
%% method: 无
%% score: 3
%% degree: 450
%% duplicateId: 0
%% body:
如图所示，示波管由电子枪竖直方向偏转电极 YY$'$、水平方向偏转电极 XX$'$ 和荧光屏组成。电极 XX$'$ 的长度为 $l$、间距为 $d$、极板间电压为 $U$，YY$'$ 极板间电压为零，电子枪加速电压为 10$U$。电子刚离开金属丝的速度为零，从电子枪射出后沿 OO$'$ 方向进入偏转电极。已知电子电荷量为 $e$，质量为 $m$，则电子（    ）
\onepicture[align=right]{figs/12.png}
%% answer: 
\jdanswer{
D
}
%% memo:
\memoanswer{
A．由牛顿第二定律可得，在 XX' 极板间的加速度大小

 $a_{x}=\frac{eE}{m}=\frac{eU}{md}$

A 错误；

B．电子电极 XX' 间运动时，有

 $v_{x}=a_{x}t$

 $t=\frac{l}{d}$

电子离开电极 XX' 时的动能为

 $E_{k}=\frac{1}{2}m(v_{0}^{2}+v_{x}^{2})=eU(10+\frac{l^{2}}{2d^{2}})$

电子离开电极 XX' 后做匀速直线运动，所以打在荧光屏时，动能大小为  $eU(10+\frac{l^{2}}{2d^{2}})$，B 错误；

C. 在 XX' 极板间受到电场力的冲量大小

 $I_{x}=mv_{x}=\frac{l}{d}\sqrt{\frac{meU}{20}}$

C 错误；

D. 打在荧光屏时，其速度方向与  $OO'$ 连线夹角  $\alpha$ 的正切

 $\tan\alpha=\frac{v_{x}}{v_{0}}=\frac{l}{20d}$

D 正确。

故选 D。
}

\item
%% number: 13
%% paperName: 2023年1月浙江选考第 13 题 3 分
%% typeId: 1
%% type: 单选题
%% chapter: 光学
%% point: 光的折射
%% method: 无
%% score: 3
%% degree: 550
%% duplicateId: 0
%% body:
如图所示为一斜边镀银的等腰直角棱镜的截面图。一细黄光束从直角边 AB 以角度 $\theta$ 入射，依次经 AC 和 BC 两次反射，从直角边 AC 出射。出射光线相对于入射光线偏转了 $\alpha$ 角，则 $\alpha$\xzanswer{A}
\onepicture{figs/13.png}
\fourchoices[answer=A]
{等于 90${^\circ}$}
{大于 90${^\circ}$}
{小于 90${^\circ}$}
{与棱镜的折射率有关}
%% answer: A
%% memo:
\memoanswer{
如图所示

\onepicture[w=6cm]{figs/13_an.jpg}

设光线在 AB 边的折射角为  $\beta$，根据折射定律可得

 $n=\frac{\sin\theta}{\sin\beta}$

设光线在 BC 边的入射角为  $\varphi$，光线在 AC 边的入射角为 r，折射角为 i；由反射定律和
几何知识可知

 $\beta+\varphi=45^{\circ}$

 $\beta+2\varphi+r=90^{\circ}$

联立解得

 $r=\beta$

根据折射定律可得

 $\frac{\sin i}{\sin r}=\frac{\sin i}{\sin \beta}=n$

可得

 $i=\theta$

过 D 点做出射光的平行线，则该平行线与 AB 的夹角为  $\theta$，由几何知识可知，入射光与出射光的夹角为  $90^{\circ}$。

故选 A。


14\. 下列说法正确的是（）

A. 利用电容传感器可制成麦克风

B．物体受合外力越大，则动量变化越快

C．利用红外传感器可制成商场的自动门

D．牛顿运动定律不适用，则动量守恒定律也不适用

【详解】A．声音使振动膜片振动，改变两极板间距离，使声音信号转换成电信号，则可以利用电容传感器可制成麦克风，故 A 正确；

B．由动量定理有  $F_{\text{合}}t = \Delta p$ 可得  $F_{\text{合}} = \frac{\Delta p}{t}$

可知，物体受合外力越大，则动量变化越快，故 B 正确；

C．人体可以向外界释放红外线，感应装置接收到红外线后，可以开门，则可以利用红外传感器可制成商场的自动门，故 C 正确；

D．牛顿运动定律只适用于宏观低速问题，不适用于微观高速问题。而动量守恒定律既适用于低速宏观问题，也适用于高速微观问题，故 D 错误。

故选 ABC。
}

\item
%% number: 14
%% paperName: 2023年1月浙江选考第 14 题 3 分
%% typeId: 2
%% type: 多选题
%% chapter: 电路
%% point: 传感器
%% method: 无
%% score: 3
%% degree: 850
%% duplicateId: 0
%% body:
下列说法正确的是\xzanswer{ABC}
\fourchoices[answer=ABC]
{利用电容传感器可制成麦克风}
{物体受合外力越大，则动量变化越快}
{利用红外传感器可制成商场的自动门}
{牛顿运动定律不适用，则动量守恒定律也不适用}
%% answer: ABC
%% memo:
\memoanswer{
本题原卷及材料中未提供详解，待补充。
}

\item
%% number: 15
%% paperName: 2023年1月浙江选考第 15 题 3 分
%% typeId: 2
%% type: 多选题
%% chapter: 光学
%% point: 光的干涉
%% method: 无
%% score: 3
%% degree: 650
%% duplicateId: 0
%% body:
氢原子从高能级向低能级跃迁时，会产生四种频率的可见光，其光谱如图 1 所示。氢原子从能级 6 跃迁到能级 2 产生可见光 I，从能级 3 跃迁到能级 2 产生可见光 Ⅱ。用同一双缝干涉装置研究两种光的干涉现象，得到如图 2 和图 3 所示的干涉条纹。用两种光分别照射如图 4 所示的实验装置，都能产生光电效应。下列说法正确的是\xzanswer{CD}
\onepicture{figs/15.png}
\fourchoices[answer=CD]
{图 1 中的 H$_{\alpha}$ 对应的是 Ⅰ}
{图 2 中的干涉条纹对应的是 Ⅱ}
{Ⅰ 的光子动量大于 Ⅱ 的光子动量}
{P 向 a 移动，电流表示数为零时 Ⅰ 对应的电压表示数比 Ⅱ 的大}
%% answer: CD
%% memo:
\memoanswer{
根据题意可知。氢原子发生能级跃迁时，由公式可得

 $E_{m}-E_{n}=h\nu=\frac{hc}{\lambda}$

可知，可见光 I 的频率大，波长小，可见光 II 的频率小，波长大。

A. 可知，图 1 中的  $H_{a}$ 对应的是可见光 II，故 A 错误；

B. 由公式有，干涉条纹间距为

 $\Delta x = \lambda \frac{l}{d}$

由图可知，图 2 中间距较小，则波长较小，对应的是可见光 I，故 B 错误；

C. 根据题意，由公式可得，光子动量为

 $p = \frac{h\nu}{c} = \frac{h}{\lambda}$

可知，I 的光子动量大于 II 的光子动量，故 C 正确；

D. 根据光电效应方程及动能定理可得

 $eU_{c}=h\nu-W_{0}$

可知，频率越大，遏止电压越大，则 P 向 a 移动，电流表示数为零时 I 对应的电压表示数比 II 的大，故 D 正确。

故选 CD。



16\. 实验题（Ⅰ、Ⅱ两题共14分）

Ⅰ．（7分）（1）在“探究小车速度随时间变化的规律”的实验中，实验装置如图1所示。

图1

① 需要的实验操作有 \_\_\_ （多选）；

A. 调节滑轮使细线与轨道平行

B. 倾斜轨道以补偿阻力

C. 小车靠近打点计时器静止释放

D. 先接通电源再释放小车

②经正确操作后打出一条纸带，截取其中一段如图 2 所示。选取连续打出的点 0、1、2、3、4 为计数点，则计数点 1 的读数为 \_\_\_ $\Ucm$。已知打点计时器所用交流电源的频率为 50 $\UHz$，则打计数点 2 时小车的速度大小为： \_\_\_ $\Ums$（结果保留 3 位有效数字）。

图2

【详解】① A. 实验需要调节滑轮使细线与轨道平行，选项 A 正确；

B. 该实验只要使得小车加速运动即可，不需要倾斜轨道补偿阻力，选项 B 错误；

C. 为了充分利用纸带，则小车靠近打点计时器静止释放，选项 C 正确；

D. 先接通电源再释放小车，选项 D 正确。

故选 ACD。

② 计数点 1 的读数为 2.75 $\Ucm$。已知打点计时器所用交流电源的频率为 50 $\UHz$，则打点周
期 T = 0.02 $\Us$，则打计数点 2 时小车的速度大小为

 $v_{2}=\frac{x_{13}}{2T}=\frac{(8.65-2.75)\times10^{-2}}{0.04}\Ums=1.48\Ums$

（2）“探究向心力大小的表达式”实验装置如图所示。

① 采用的实验方法是 \_\_\_

A. 控制变量法 B. 等效法 C. 模拟法

② 在小球质量和转动半径相同的情况下，逐渐加速转动手柄到一定速度后保持匀速转动。此时左右标尺露出的红白相间等分标记的比值等于两小球的\_\_\_之比（选填“线速度大小”、“角速度平方”或“周期平方”）；在加速转动手柄过程中，左右标尺露出红白相何等分标记的比值\_\_\_（选填“不变”、“变大”或“变小”）。

【详解】① 本实验先控制住其它几个因素不变，集中研究其中一个因素变化所产生的影响，采用的实验方法是控制变量法；故选 A。

② 标尺上露出的红白相间的等分格数之比为两个小球所受向心力的比值，根据  $F = m r \omega^{2}$ 可知比值等于两小球的角速度平方之比。

小球质量和转动半径相同的情况下，由于两球线速度相同，根据  $\omega = \frac{V}{R}$ 可知两球角速度相同；又根据  $F = m r \omega^{2}$ 可知向心力之比为 1:1，逐渐加大手柄的转速，左右标尺露出的格数之比不变。

Ⅱ.（7分）在“测量金属丝的电阻率”实验中：

（1）测量一段金属丝电阻时所用器材和部分电路连线如图 1 所示，图中的导线 a 端应与 \_\_\_（选填 “-”、“0.6” 或 “3”）接线柱连接，b 端应与 \_\_\_（选填 “-”、“0.6” 或 “3”）接线柱连接。开关闭合前，图 1 中滑动变阻器滑片应置于 \_\_\_（选填 “左” 或 “右”）端。

图1

（2）合上开关，调节滑动变阻器，得到多组 U 和 I 数据。甲同学由每组 U、I 数据计算电阻，然后求电阻平均值；乙同学通过 U-I 图像求电阻。则两种求电阻的方法更合理的是 \_\_\_（选填 “甲” 或 “乙”）。

（3）两同学进一步探究用镍铬丝将满偏电流  $I_{g}=300\ \mu A$ 的表头 G 改装成电流表。如图 2 所示，表头 G 两端并联长为 L 的镍铬丝，调节滑动变阻器使表头 G 满偏，毫安表示数为
I。改变 L，重复上述步骤，获得多组 I、L 数据，作出  $I - \frac{1}{L}$ 图像如图 3 所示。

图2

图3

则  $I - \frac{1}{L}$ 图像斜率  $k =$ \_\_\_ $\UmA\cdot\Um$。若要把该表头 G 改装成量程为 9 $\UmA$ 的电流表，需要把长为 \_\_\_ $\Um$ 的镍铬丝并联在表头 G 两端。（结果均保留两位有效数字）

【详解】（1）实验中用两节干电池供电，则图中的导线 a 端应与 “0.6” 接线柱连接，电压表测电阻两端的电压，则金属丝的电阻较小，电流表外接误差较小，故 b 端应与 “0.6” 接线柱连接。为了保护电表，开关闭合前，图 1 中滑动变阻器滑片应置于左端。

（2）乙同学通过 U-I 图像求电阻，可减少实验的误差，则乙同学求电阻的方法更合理；

（3）由图像可知  $I - \frac{1}{L}$ 图像斜率

 $k=\frac{9.0-4.8}{3.8-2.0}=2.3\UmA\cdot\Um$

方法一：由电路可知

 $I_{g}r_{g}=(I-I_{g})\rho\frac{L}{S}$

解得

 $I=I_{g}+\frac{I_{g}r_{g}S}{\rho}\cdot\frac{1}{L}$

 $I_{g}=0.3\UmA$

由图像，则

 $\frac{I_{g}r_{g}S}{\rho}=2.3\UmA\cdot\Um$

若要把该表头 G 改装成量程为 9 $\UmA$ 的电流表，则并联的电阻
$$ \rho\frac{L}{S}=\frac{I_{g}r_{g}}{I^{\prime}-I_{g}} $$

即

$$ \frac{\rho}{I_{g}r_{g}S}L^{\prime}=\frac{1}{I^{\prime}-I_{g}} $$

解得  $L' = 0.26\Um$

方法二：延长图像可知，当  $I = 9.0\UmA$ 时可得  $\frac{1}{L} = 3.8$，即  $L' = 0.26\Um$
}

\item
%% number: 16
%% paperName: 2023年1月浙江选考第 16 题 4 分
%% typeId: 4
%% type: 实验题
%% chapter: 直线运动
%% point: 打点计时器公式
%% method: 无
%% score: 4
%% degree: 750
%% duplicateId: 0
%% body:
在“探究小车速度随时间变化的规律”的实验中，实验装置如图 1 所示。
\begin{enumerate}
	\item
	需要的实验操作有\_\_\_\_\_\_\_\_\_\_（多选）；
	\item
	经正确操作后打出一条纸带，截取其中一段如图 2 所示。选取连续打出的点 0、1、2、3、4 为计数点，则计数点 1 的读数为\_\_\_\_\_\_\_\_\_\_ $\Ucm$。已知打点计时器所用交流电源的频率为 50 $\UHz$，则打计数点 2 时小车的速度大小为：\_\_\_\_\_\_\_\_\_\_$\Ums$（结果保留 3 位有效数字）。
\end{enumerate}
\twopicture[labelA=20231月浙江16a, labelB=20231月浙江16b, align=right]
{figs/16a.png}
{figs/16b.png}
%% answer: 
\jdanswer{
\begin{enumerate}
	\item
	ACD

	\item
	2.75，1.48

\end{enumerate}
}
%% memo:
\memoanswer{
（1）A．实验需要调节滑轮使细线与轨道平行，选项 A 正确；

B．该实验只要使得小车加速运动即可，不需要倾斜轨道补偿阻力，选项 B 错误；

C．为了充分利用纸带，则小车靠近打点计时器静止释放，选项 C 正确；

D．先接通电源再释放小车，选项 D 正确。故选 ACD。

（2）计数点 1 的读数为 $2.75\Ucm$。打点周期 $T = 0.02\Us$，则打计数点 2 时小车的速度大小为 $v_{2}=\frac{x_{13}}{2T}=\frac{(8.65-2.75)\times10^{-2}}{0.04}\Ums=1.48\Ums$。
}

\item
%% number: 16
%% paperName: 2023年1月浙江选考第 16 题 3 分
%% typeId: 4
%% type: 实验题
%% chapter: 曲线运动
%% point: 探究向心力实验
%% method: 无
%% score: 3
%% degree: 750
%% duplicateId: 0
%% body:
“探究向心力大小的表达式”实验装置如图所示。
\begin{enumerate}
	\item
	采用的实验方法是\_\_\_\_\_\_\_\_\_\_
	\item
	在小球质量和转动半径相同的情况下，逐渐加速转动手柄到一定速度后保持匀速转动。此时左右标尺露出的红白相间等分标记的比值等于两小球的\_\_\_\_\_\_\_\_之比（选填“线速度大小”、“角速度平方”或“周期平方”）；在加速转动手柄过程中，左右标尺露出红白相何等分标记的比值\_\_\_\_\_\_\_\_\_\_（选填“不变”、“变大”或“变小”）。
\end{enumerate}
\onepicture[align=right]{figs/16.png}
%% answer: 
\jdanswer{
\begin{enumerate}
	\item
	A

	\item
	角速度平方，不变

\end{enumerate}
}
%% memo:
\memoanswer{
（1）本实验先控制住其它几个因素不变，集中研究其中一个因素变化所产生的影响，采用的实验方法是控制变量法，故选 A。

（2）标尺上露出的红白相间的等分格数之比为两个小球所受向心力的比值，根据 $F = mr\omega^{2}$ 可知比值等于两小球的角速度平方之比。小球质量和转动半径相同的情况下，由于两球线速度相同，根据 $\omega = \frac{v}{R}$ 可知两球角速度相同；又根据 $F = mr\omega^{2}$ 可知向心力之比为 1:1，逐渐加大手柄的转速，左右标尺露出的格数之比不变。
}

\item
%% number: 16
%% paperName: 2023年1月浙江选考第 16 题 7 分
%% typeId: 4
%% type: 实验题
%% chapter: 电路
%% point: 测量电阻率
%% method: 无
%% score: 7
%% degree: 550
%% duplicateId: 0
%% body:
在“测量金属丝的电阻率”实验中：
\begin{enumerate}
	\item
	测量一段金属丝电阻时所用器材和部分电路连线如图 1 所示，图中的导线 a 端应与\_\_\_\_\_\_\_\_（选填“$-$”、“0.6”或“3”）接线柱连接，b 端应与\_\_\_\_\_\_\_\_\_\_\_（选填“$-$”、“0.6”或“3”）接线柱连接。开关闭合前，图 1 中滑动变阻器滑片应置于\_\_\_\_\_\_\_\_\_\_\_\_（选填“左”或“右”）端。
	\item
	合上开关，调节滑动变阻器，得到多组 $U$ 和 $I$ 数据。甲同学由每组 $U$、$I$ 数据计算电阻，然后求电阻平均值；乙同学通过 $U - I$ 图像求电阻。则两种求电阻的方法更合理的是\_\_\_\_\_\_\_\_\_\_\_（选填“甲”或“乙”）。
	\item
	两同学进一步探究用镍铬丝将满偏电流 $I_{g}$ = 300 $\mu$A 的表头 G 改装成电流表。如图 2 所示，表头 G 两端并联长为 $L$ 的镍铬丝，调节滑动变阻器使表头 G 满偏，毫安表示数为 $I$。改变 $L$，重复上述步骤，获得多组 $I$、$L$ 数据，作出 $I - \frac{1}{L}$ 图像如图 3 所示。
\end{enumerate}
\twopicture[labelA=20231月浙江16d, labelB=20231月浙江16e, align=right]
{figs/16d.png}
{figs/16e.png}
%% answer: 
\jdanswer{
\begin{enumerate}
	\item
	0.6，0.6，左

	\item
	乙

	\item
	2.3，0.26

\end{enumerate}
}
%% memo:
\memoanswer{
（1）实验中用两节干电池供电，则图中的导线 a 端应与“0.6”接线柱连接；电压表测电阻两端的电压，金属丝的电阻较小，电流表外接误差较小，故 b 端应与“0.6”接线柱连接。为了保护电表，开关闭合前，图 1 中滑动变阻器滑片应置于左端。

（2）乙同学通过 $U-I$ 图像求电阻，可减少实验的误差，则乙同学求电阻的方法更合理。

（3）由图像可知 $I - \frac{1}{L}$ 图像斜率 $k=\frac{9.0-4.8}{3.8-2.0}=2.3\UmA\cdot\Um$。由电路可知 $I_{g}r_{g}=(I-I_{g})\rho\frac{L}{S}$，解得 $I=I_{g}+\frac{I_{g}r_{g}S}{\rho}\cdot\frac{1}{L}$，$I_{g}=0.3\UmA$。由图像，则 $\frac{I_{g}r_{g}S}{\rho}=2.3\UmA\cdot\Um$。若要把该表头 G 改装成量程为 9 $\UmA$ 的电流表，则并联的电阻 $\rho\frac{L}{S}=\frac{I_{g}r_{g}}{I^{\prime}-I_{g}}$，即 $\frac{\rho}{I_{g}r_{g}S}L^{\prime}=\frac{1}{I^{\prime}-I_{g}}$，解得 $L' = 0.26\Um$。
}

\item
%% number: 17
%% paperName: 2023年1月浙江选考第 17 题 8 分
%% typeId: 6
%% type: 计算题
%% chapter: 气体性质
%% point: 气体多过程
%% method: 无
%% score: 8
%% degree: 550
%% duplicateId: 0
%% body:
某探究小组设计了一个报警装置，其原理如图所示。在竖直放置的圆柱形容器内用面积 $S$ = 100 $\Ucmq$、质量 $m$ = 1 $\Ukg$ 的活塞密封一定质量的理想气体，活塞能无摩擦滑动。开始时气体处于温度 $T_{A}$ = 300 $\UK$、活塞与容器底的距离 $h_{0}$ = 30 $\Ucm$ 的状态 A。环境温度升高时容器内气体被加热，活塞缓慢上升 $d$ = 3 $\Ucm$ 恰好到达容器内的卡口处，此时气体达到状态 B。活塞保持不动，气体被继续加热至温度 $T_{C}$ = 363 $\UK$ 的状态 C 时触动报警器。从状态 A 到状态 C 的过程中气体内能增加了 $\Delta U$ = 158 $\UJ$。取大气压 $p_{0}$ = 0.99$\times$10$^{5}$ $\UPa$，求气体
\begin{enumerate}
	\item
	在状态 B 的温度；
	\item
	在状态 C 的压强；
	\item
	由状态 A 到状态 C 过程中从外界吸收热量 $Q$。
\end{enumerate}
\onepicture[align=right]{figs/17.png}
%% answer: 
\jdanswer{
\begin{enumerate}
	\item
	330 $\UK$

	\item
	1.09$\times$10$^{5}$ $\UPa$

	\item
	187.7 $\UJ$

\end{enumerate}
}
%% memo:
\memoanswer{
（1）根据题意可知，气体由状态 A 变化到状态 B 的过程中，封闭气体的压强不变为  $p_{0}$，则有

 $\frac{V_{A}}{T_{A}}=\frac{V_{B}}{T_{B}}$

解得

 $T_{B}=\frac{V_{B}}{V_{A}}T_{A}=330\UK$

（2）根据题意可知，气体由状态 B 变化到状态 C 的过程中，气体的体积不变，则有  $\frac{p_{B}}{T_{B}} = \frac{p_{C}}{T_{C}}$

解得

 $p_{C}=\frac{T_{C}}{T_{B}}p_{B}=1.1p_{0}\approx1.09\times10^{5}Pa$

（3）根据题意可知，从状态 A 到状态 C 的过程中气体对外做功为  $W_{0}=p_{0}\Delta V=29.7\UJ$

由热力学第一定律有

 $\Delta U = W + Q$

解得  $Q = \Delta U + W = 187.7\UJ$
}

\item
%% number: 18
%% paperName: 2023年1月浙江选考第 18 题 11 分
%% typeId: 6
%% type: 计算题
%% chapter: 运动和力
%% point: 动量守恒定律
%% method: 无
%% score: 11
%% degree: 450
%% duplicateId: 0
%% body:
一游戏装置竖直截面如图所示，该装置由固定在水平地面上倾角 $\theta$ = 37${^\circ}$ 的直轨道 AB、螺旋圆形轨道 BCDE，倾角 $\theta$ = 37${^\circ}$ 的直轨道 EF、水平直轨道 FG 组成，除 FG 段外各段轨道均光滑，且各处平滑连接。螺旋圆形轨道与轨道 AB、EF 相切于 B(E) 处。凹槽 GHIJ 底面 HI 水平光滑，上面放有一无动力摆渡车，并紧靠在竖直侧壁 GH 处，摆渡车上表面与直轨道下 FG、平台 JK 位于同一水平面。已知螺旋圆形轨道半径 $R$ = 0.5 $\Um$，B 点高度为 1.2$R$，FG 长度 $L_{FG}$ = 2.5 $\Um$，HI 长度 $L_{0}$ = 9 $\Um$，摆渡车长度 $L$ = 3 $\Um$、质量 $m$ = 1 $\Ukg$。将一质量也为 $m$ 的滑块从倾斜轨道 AB 上高度 $h$ = 2.3 $\Um$ 处静止释放，滑块在 FG 段运动时的阻力为其重力的 0.2 倍。（摆渡车碰到竖直侧壁 IJ 立即静止，滑块视为质点，不计空气阻力，sin37${^\circ}$ = 0.6，cos37${^\circ}$ = 0.8）
\begin{enumerate}
	\item
	求滑块过 C 点的速度大小 $v_{C}$ 和轨道对滑块的作用力大小 $F_{C}$；
	\item
	摆渡车碰到 IJ 前，滑块恰好不脱离摆渡车，求滑块与摆渡车之间的动摩擦因数 $\mu$；
	\item
	在（2）的条件下，求滑块从 G 到 J 所用的时间 $t$。
\end{enumerate}
\onepicture[align=right]{\includesvg{figs/18}}
%% answer: 
\jdanswer{
\begin{enumerate}
	\item
	$v_{C}$ = 4 $\Ums$，$F_{C}$ = 22 $\UN$

	\item
	$\mu$ = 0.3

	\item
	$t$ = 2.5 $\Us$

\end{enumerate}
}
%% memo:
\memoanswer{
（1）滑块从静止释放到 C 点过程，根据动能定理可得

 $mg(h-1.2R-R-R\cos\theta)=\frac{1}{2}mv_{C}^{2}$

解得  $v_{C}=4\ \Ums$

滑块过 C 点时，根据牛顿第二定律可得

 $F_{C}+mg=m\frac{v_{C}^{2}}{R}$

解得  $F_{C}=22\ N$

（2）设滑块刚滑上摆渡车时的速度大小为 v，从静止释放到 G 点过程，根据动能定理可得

 $mgh - 0.2mgL_{FG} = \frac{1}{2}mv^{2}$

解得  $v = 6 \, \Ums$

摆渡车碰到 IJ 前，滑块恰好不脱离摆渡车，说明滑块到达摆渡车右端时刚好与摆渡车共速  $v_{1}$，以滑块和摆渡车为系统，根据系统动量守恒可得

 $mv = 2mv_{1}$

解得  $v_{1}=\frac{v}{2}=3\Ums$

根据能量守恒可得

 $Q=\mu mgL=\frac{1}{2}mv^{2}-\frac{1}{2}\times2mv_{1}^{2}$

解得  $\mu = 0.3$

（3）滑块从滑上摆渡车到与摆渡车共速过程，滑块的加速度大小为

 $a=\frac{\mu mg}{m}=3\Umsq$

所用时间为

 $t_{1}=\frac{v-v_{1}}{a}=1\Us$

此过程滑块通过的位移为

 $x_{1}=\frac{v+v_{1}}{2}t_{1}=4.5\Um$

滑块与摆渡车共速后，滑块与摆渡车一起做匀速直线运动，该过程所用时间为
$$ t_{2}=\frac{L_{0}-x_{1}}{v_{1}}=1.5\Us $$

则滑块从 G 到 J 所用的时间为

 $t = t_{1} + t_{2} = 2.5\Us$
}

\item
%% number: 19
%% paperName: 2023年1月浙江选考第 19 题 11 分
%% typeId: 6
%% type: 计算题
%% chapter: 电磁感应
%% point: 法拉第电磁感应定律
%% method: 无
%% score: 11
%% degree: 450
%% duplicateId: 0
%% body:
如图 1 所示，刚性导体线框由长为 $L$、质量均为 $m$ 的两根竖杆，与长为 2$l$ 的两轻质横杆组成，且 $L \gg$ 2$l$。线框通有恒定电流 $I_{0}$，可以绕其中心竖直轴转动。以线框中心 O 为原点、转轴为 $z$ 轴建立直角坐标系，在 $y$ 轴上距离 $O$ 为 $a$ 处，固定放置一半径远小于 $a$，面积为 $S$、电阻为 $R$ 的小圆环，其平面垂直于 $y$ 轴。在外力作用下，通电线框绕转轴以角速度 $\omega$ 匀速转动，当线框平面与 $xOy$ 平面重合时为计时零点，圆环处的磁感应强度的 $y$ 分量 $B_{y}$ 与时间的近似关系如图 2 所示，图中 $B_{0}$ 已知。
\begin{enumerate}
	\item
	求 0 到 $\frac{\pi }{\omega }$ 时间内，流过圆环横截面的电荷量 $q$；
	\item
	沿 $y$ 轴正方向看以逆时针为电流正方向，在 0 \textasciitilde{} $\frac{{2\pi }}{{3\omega }}$ 时间内，求圆环中的电流与时间的关系；
	\item
	求圆环中电流的有效值；
	\item
	当撤去外力，线框将缓慢减速，经 $\frac{\pi }{\omega }$ 时间角速度减小量为 $\Delta\omega$（$\frac{{\Delta \omega }}{\omega } \ll$ 1），设线框与圆环的能量转换效率为 $k$，求 $\Delta\omega$ 的值（当 0 < $x \ll$ 1，有 (1 $- x$)$^{2} \approx$ 1 - 2$x$。
\end{enumerate}
\onepicture[align=right]{figs/19.png}
%% answer: 
\jdanswer{
\begin{enumerate}
	\item
	$\frac{{2{B_0}S}}{R}$

	\item
	当 0 $\leq t \leq \frac{\pi }{{3\omega }}$，$I$ = 0

	\item
	当  $\frac{\pi }{{3\omega }} < t \leq \frac{{2\pi }}{{3\omega }}$，$I = \frac{{6\omega {B_0}S}}{{\pi R}}$

	\item
	$\frac{{2\sqrt 3 \omega {B_0}S}}{{\pi R}}$

	\item
	$\frac{{6B_0^2{S^2}}}{{\pi km{l^2}R}}$

\end{enumerate}
}
%% memo:
\memoanswer{
（1）由法拉第电磁感应定律

 $E=\frac{\Delta\Phi}{\Delta t}$

由闭合电路欧姆定律

 $I=\frac{E}{R}$

由电流定义式

 $q=\sum I\Delta t$

联立可得

 $q=\frac{\Delta\Phi}{R}=\frac{2B_{0}S}{R}\quad q=\frac{2B_{0}S}{R}$
(2) 在  $0 \sim \frac{\pi}{3\omega}$ 时

 $I=\frac{E}{R}=\frac{\Delta\Phi}{R\Delta t}=0$

在  $\frac{\pi}{3\omega} \sim \frac{2\pi}{3\omega}$ 时

 $I=\frac{E}{R}=\frac{\frac{\Delta B}{\Delta t}S}{R}=-\frac{6\omega B_{0}S}{\pi R}$

（3）从能量角度

 $I^{2}R\frac{\pi}{3\omega}=I_{\text{有}}^{2}R\frac{\pi}{\omega}$

解得

 $I_{\text{有}} = \frac{2\sqrt{3\omega B_{0}S}}{\pi R}$

（4）由能量传递

 $k\cdot\frac{1}{2}\cdot2m\left[\omega^{2}l^{2}-\left(\omega-\Delta\omega\right)^{2}l^{2}\right]=I_{\text{有}}^{2}R\frac{\pi}{\omega}$

化简可得

 $km\omega^{2}l^{2}\left[1-\left(1-\frac{\Delta\omega}{\omega}\right)^{2}\right]=\frac{12\omega^{2}B_{0}^{2}S^{2}}{\pi^{2}R^{2}}\cdot R\cdot\frac{\pi}{\omega}$

即

 $km\omega^{2}l^{2}\cdot2\frac{\Delta\omega}{\omega}=\frac{12\omega^{2}B_{0}^{2}S^{2}}{\pi R\omega}$

解得  $\Delta\omega = \frac{6B_{0}^{2}S^{2}}{\pi km l^{2}R}$
}

\item
%% number: 20
%% paperName: 2023年1月浙江选考第 20 题 11 分
%% typeId: 6
%% type: 计算题
%% chapter: 磁场
%% point: 带电粒子在组合场中的运动
%% method: 无
%% score: 11
%% degree: 450
%% duplicateId: 0
%% body:
探究离子源发射速度大小和方向分布的原理如图所示。$x$ 轴上方存在垂直 $xOy$ 平面向外、磁感应强度大小为 $B $的匀强磁场。$x$ 轴下方的分析器由两块相距为 $d$、长度足够的平行金属薄板 M 和 N 组成，其中位于 $x$ 轴的 M 板中心有一小孔 C（孔径忽略不计），N 板连接电流表后接地。位于坐标原点 $O$ 的离子源能发射质量为 $m$、电荷量为 $q$ 的正离子，其速度方向与 $y$ 轴夹角最大值为 60${^\circ}$；且各个方向均有速度大小连续分布在 $\frac{1}{2} v_{0}$ 和 $\sqrt 2  v_{0}$ 之间的离子射出。已知速度大小为 $v_{0}$、沿 $y$ 轴正方向射出的离子经磁场偏转后恰好垂直 $x$ 轴射入孔 C。未能射入孔 C 的其它离子被分析器的接地外罩屏蔽（图中没有画出）。不计离子的重力及相互作用，不考虑离子间的碰撞。
\begin{enumerate}
	\item
	求孔 C 所处位置的坐标 $x_{0}$；
	\item
	求离子打在 N 板上区域的长度 $L$；
	\item
	若在 N 与 M 板之间加载电压，调节其大小，求电流表示数刚为 0 时的电压 $U_{0}$；
	\item
	若将分析器沿着 $x$ 轴平移，调节加载在 N 与 M 板之间的电压，求电流表示数刚为 0 时的电压 $U_{x}$ 与孔 C 位置坐标 $x$ 之间关系式。
\end{enumerate}
\onepicture[align=right]{figs/20.png}
%% answer: 
\jdanswer{
\begin{enumerate}
	\item
	$\frac{{2m{v_0}}}{{qB}}$

	\item
	2$d$

	\item
	$\frac{{mv_0^2}}{{2q}}$

	\item
	当 $\frac{{m{v_0}}}{{2qB}} < x < \frac{{2\sqrt 2 m{v_0}}}{{qB}}$ 时，$U_{x} = \frac{{q{B^2}{x^2}}}{{8m}}$

\end{enumerate}
}
%% memo:
\memoanswer{
速度大小为 $v_{0}$、沿 $y$ 轴正方向射出的离子经磁场偏转后轨迹如图

\onepicture[w=6cm]{figs/20_an.png}

$\ldots\ldots\ldots\ldots\ldots\ldots\ldots\ldots\ldots\ldots$
}

\end{enumerate}

\end{document}
